\chapter{笼空间中的解析的高阶柯西坐标}

% ===== Chapter 5 macros =====
\providecommand{\Bezier}{B\'{e}zier }
\providecommand{\tr}[1]{#1}
\providecommand{\tb}[1]{#1}
\providecommand{\todo}[1]{}
\providecommand{\argmin}{\mathop{\mathrm{arg\,min}}\limits}
\providecommand{\argmax}{\mathop{\mathrm{arg\,max}}\limits}
\providecommand{\res}{\mathop{\mathrm{Res}}}
\providecommand{\re}{\mathop{\mathrm{Re}}}
\providecommand{\im}{\mathop{\mathrm{Im}}}

\makeatletter
\@ifundefined{theorem}{\newtheorem{theorem}{Theorem}}{}
\@ifundefined{lemma}{\newtheorem{lemma}{Lemma}}{}
\@ifundefined{prop}{%
  \@ifundefined{theorem}{%
    \newtheorem{prop}{Proposition}[chapter]%
  }{%
    \newtheorem{prop}[theorem]{Proposition}%
  }%
}{}
\makeatother

\def \ma {\mathbf{a}}
\def \me {\mathbf{e}}
\def \mb {\mathbf{b}}
\def \mc {\mathbf{c}}
\def \md {\mathbf{d}}
\def \mf {\mathbf{f}}
\def \mg {\mathbf{g}}
\def \mh {\mathbf{h}}
\def \mi {\mathbf{i}}
\def \mj {\mathbf{j}}
\def \mk {\mathbf{k}}
\def \ml {\mathbf{l}}
\def \mm {\mathbf{m}}
\def \mn {\mathbf{n}}
\def \mo {\mathbf{o}}
\def \mp {\mathbf{p}}
\def \mq {\mathbf{q}}
\def \mr {\mathbf{r}}
\def \ms {\mathbf{s}}
\def \mt {\mathbf{t}}
\def \muu {\mathbf{u}}
\def \mv {\mathbf{v}}
\def \mw {\mathbf{w}}
\def \mx {\mathbf{x}}
\def \my {\mathbf{y}}
\def \mz {\mathbf{z}}

\section{概述}\label{sec:chap5-cauchy-overview}
本章从笼空间中的映射问题出发。与壳空间映射不同，笼空间映射以笼边界围成的控制域为定义域，需要将边界上的控制量稳定传播到内部区域，因此其关键在于构造兼具插值性、线性精确性与良好形变质量的坐标函数。笼坐标为这一类问题提供了统一表述：用户只需在边界上给定稀疏约束，即可将其传播到整个内部区域。由于具有直观的边界控制能力，这类方法已被广泛用于二维形状变形、图像着色和艺术编辑等任务。现有笼坐标大致可分为两类：一类通过边界顶点值的加权组合实现插值，如均值坐标、正均值坐标与调和坐标~\cite{Floater2003,Hormann2006,Lipman2007,Joshi2007}；另一类进一步引入边界导数信息，如 Green 坐标与 Cauchy 坐标~\cite{Lipman2008,Weber2009}。其中，Green 坐标与 Cauchy 坐标在二维情形下都与保角形变密切相关，因此在形状保持方面具有明显优势。

\begin{figure}[t]
	\centering
	\begin{overpic}[width=0.99\linewidth]{figures/chap4_teaser.pdf}
		{
			\put(10,20){\small 输入}
			\put(75,21){\small 输入}
			\put(15,2){\small 变形后}
			\put(70,-2){\small 变形后}
		}
	\end{overpic}
	\vspace{4mm}
	\caption{
		利用本章坐标及其导数进行变形。上排为输入图像与三次笼，下排分别为基于笼的变形与点到点变形。
	}
	\label{fig:cau-teaser}
\end{figure}

当笼空间映射面对具有明显曲边特征的输入形状时，线性笼往往需要大量边段才能贴合轮廓，这不仅削弱了交互操控的直观性，也容易影响边界附近的形变光滑性。近年来，三次均值坐标、多项式 Green 坐标和多项式 Cauchy 坐标等工作开始支持将初始直边笼推广为高阶多项式曲线~\cite{Li2013,MichelThiery2023,Lin2024PolynomialCC}。沿着这一思路，本章关注由多项式曲线组成的输入笼，并将其称为高阶笼。与线性笼相比，高阶笼能够更紧密地逼近原始形状，以更少的参数表达边界；同时，多项式边界与 CAD 中常见的自由形状设计方式更一致，更适合作为直接的编辑界面；此外，高阶笼诱导的变形在边界附近也更光滑。

\begin{figure}[t]
	\centering
	\begin{overpic}[width=0.95\linewidth]{figures/chap4_compare_cagenew.pdf}
		{
			\put(46,68){\small 线性笼}
			\put(47,63){\small 72 段}
			\put(46,47){\small 线性笼}
			\put(46,41.5){\small 120 段}
			\put(46,24){\small 三次笼}
			\put(47,18){\small 24 段}
		}
	\end{overpic}
	\vspace{2mm}
	\caption{
		使用线性输入笼与高阶输入笼的变形对比。在线性笼下，需要大量边段才能包围复杂曲边形状；而高阶输入笼只需更少控制量即可产生更自然的变形。
	}
	\label{fig:cau-cmp-cages}
\end{figure}

在笼空间映射中，Cauchy 坐标尤其适用于强调角度保持与光滑传播的编辑任务。尽管已有工作将 Cauchy 坐标推广到高阶笼~\cite{Lin2024PolynomialCC}，但其仍需借助中间线性笼，并通过数值方式计算从中间笼到输入笼映射的逆映射。这带来两个问题：其一，坐标及其导数缺乏闭式表达，难以直接支持点到点变形；其二，高阶输入笼需要满足较强且不够直观的几何约束，因此难以处理任意有效的高阶包围边界。围绕这一问题，本章给出二维高阶笼上的 Cauchy 坐标及其导数的闭式形式。核心在于将相关积分转化为有理多项式积分，并结合对数函数与留数定理构造可计算的解析表达；同一积分框架还可扩展到高阶笼上的 Green 坐标。

接下来，\S~\ref{sec:cau-method} 节给出高阶笼上的 Cauchy 坐标推导、闭式解、扩展与实现细节，\S~\ref{sec:cau-results} 节展示其在基于笼变形和点到点变形中的效果与对比实验。

\section{高阶柯西坐标} \label{sec:cau-method}
本章先在第~\ref{sec:cau-revisit} 节回顾 Cauchy 坐标，再在第~\ref{sec:cau-High-order-Cauchy} 节介绍笼上的 高阶Cauchy 坐标，最后在第~\ref{sec:cau-closed-form} 节推导这些坐标及其导数的闭式表达。
%
更多讨论见第~\ref{sec:cau-connection} 与第~\ref{sec:cau-details} 节。

\subsection{ Cauchy 坐标}\label{sec:cau-revisit}
作为复分析中的基础结论，Cauchy 积分公式指出：定义在平面区域 \(\Omega\) 上、且边界 $\partial \Omega$ 为简单闭曲线的全纯函数，可由其在 $\partial \Omega$ 上的取值完全决定。
%
将连续函数 \(f\) 的边界值代入 Cauchy 积分公式右侧，可得到一个全纯函数 \(u\)：
\begin{equation}\label{equ:cau-Cauchy-integral}
    u(z)=\frac{1}{2\pi i}\int_{\partial \Omega}\frac{f(\zeta)}{\zeta-z} d\zeta,\ \ z \in \Omega,
\end{equation}
其中 $z$ 为复数，且 \(i=\sqrt{-1}\)。
这种从 \(f\) 到 \(u\) 的变换称为\emph{Cauchy 变换} \cite{Bell2015cauchytransform}。
文献中的两类 Cauchy 坐标~\cite{Weber2009,Lin2024PolynomialCC} 均来源于该变换。

\paragraph{笼（Cage）}
给定有界区域 $\Omega$，其边界 $\partial \Omega$（即\emph{笼}）是一个无自交闭合曲线，由顶点集 $V = \{\mv_j\}^n_{j=1}$ 和边集 $E = \{\me_j\}^n_{j=1}$ 组成（按逆时针定向），见图~\ref{fig:cau-notations}。
%
在复平面上，每个顶点 $\mv_j$ 记作复数 \(z_j=\mv_{j,x}+\mv_{j,y} i\)。
%
Cauchy 坐标~\cite{Weber2009} 与多项式 Cauchy 坐标~\cite{Lin2024PolynomialCC} 都定义在具有直边的初始笼上，即对边 $\me_j = \overline{\mv_{j-1}\mv_j}$，有 \( z_{\me_j}(t) = (1-t) z_{j-1} + t z_j,\ t\in[0,1]\)。
二者的区别在于变形后笼边的表示：
\begin{itemize}
    \item 对 Cauchy 坐标，新的 $\me_j$ 为线段，即 $f^\text{new}(z_{\me_j}(t)) = (1-t) p^\text{new}_{j-1}+t p^\text{new}_{j}$，其中 $p^\text{new}_{j}$ 为顶点 $\mv_j$ 的新位置。
    \item 对多项式 Cauchy 坐标，新的 $\me_j$ 为次数 $n_j \geq 1$ 的多项式曲线，即 $f^\text{new}(z_{\me_j}(t)) = \sum_{k=0}^{n_j} B^{n_j}_k(t) p^\text{new}_{\me_j,k}$，其中 $B^{n_j}_k (t)$ 为 Bernstein 多项式，$p^\text{new}_{\me_j,k} \in \mathbb{C}$ 为控制点。
\end{itemize}
由于多项式情形包含线性情形（$n_j = 1$），下文仅讨论多项式 Cauchy 坐标。

\paragraph{Cauchy 坐标推导}
将变形笼的新位置代入式~\eqref{equ:cau-Cauchy-integral}，可得
\begin{equation}\label{equ:cau-discretization}
\begin{split}
    u(z) &:=\frac{1}{2\pi i} \sum_{j=1}^n \int_{\me_j}\frac{f^\text{new}(\zeta)}{\zeta-z}d\zeta \\
     &=\frac{1}{2 \pi i}\sum_{j=1}^n \sum_{k=0}^{n_j} \left(\int_{\me_j}\frac{B^{n_j}_k (t)}{\zeta-z}d\zeta\right) p^\text{new}_{\me_j,k}, \ \ z \in \Omega.
\end{split}
\end{equation}
用户通过编辑控制点 $p^\text{new}_{\me_j,k}$ 可直观地完成形变。
此外，在式~\eqref{equ:cau-discretization} 中，与控制点 $p^\text{new}_{\me_j,k}$ 对应的全部积分项（$\frac{1}{2 \pi i}\left(\int_{\me_j}\frac{B^{n_j}_k (t)}{\zeta-z}d\zeta\right)$）共同构成其 Cauchy 坐标~\cite{Lin2024PolynomialCC}。
由于 $B^{n_j}_k (t) = 
\begin{pmatrix}
n_j \\
k \\
\end{pmatrix}
t^k (1-t)^{n_j-k}$，
只需评估不同 $k$ 下的积分 $\left(\int_{\me_j}\frac{t^k}{\zeta-z}d\zeta\right)$ 即可得到坐标。

\begin{figure}[t]
  \centering
   \vspace{-2mm}
  \begin{overpic}[width=0.99\linewidth]{chap4_notation}
    {
 \put(31,9){\small $\me_j$}
 \put(40,19){\small $\zeta$}
 \put(13,13.5){\small $\Omega$}
 \put(23,23.5){\small $z$}
 \put(50,24){\small $u$}
 \put(80,23.5){\small $u(z)$}
 \put(91,7){\small $f^{\text{new}}(\me_j)$}
    }
  \end{overpic}
  \vspace{-5mm}
  \caption{
  高阶笼变形示意。原始曲线 $\me_j$ 变形为 $f^{\text{new}(\me_j)}$，笼内任意点 $z$ 变形为 $u(z)$。
  }
  \label{fig:cau-notations}
\end{figure}

\subsection{高阶笼上的二维 Cauchy 坐标及其导数}\label{sec:cau-High-order-Cauchy}

\paragraph{更强的变形交互性}
由 Cauchy 积分公式可得
\begin{equation}
    u(z)-z=\frac{1}{2\pi i}\int_{\partial \Omega}\frac{f(\zeta)-\zeta}{\zeta-z} d\zeta,\ \ z \in \Omega,
\end{equation}
其中 \(f(\zeta)-\zeta\) 表示笼边界点 $\zeta$ 的位移（通常非全纯），而 \(u(z)-z\) 为内部点 \(z\) 的形变位移。
由于 \(u(z)-z\) 是全纯函数，依据最大模原理，其最大位移出现在 \(\partial \Omega\) 上。
因此，对同一边界位移场而言，离笼越近的区域形变位移越大，边界附近交互性更强。
这也是为何在交互编辑中，通常希望笼尽可能贴近形状，而非采用远离形状的大包围边界。
高阶笼相比线性笼具有更强逼近能力，因此我们将 Cauchy 坐标推广到高阶输入笼。

\paragraph{边界上的全纯性}
Cauchy 坐标与多项式 Cauchy 坐标都能在 \(\Omega\) 内生成全纯函数 $u$。
若全纯函数导数不为 0，则其为保角映射，因此可得到保角形变。
但 $u$ 在 \(\partial \Omega\) 上一般并非全纯，这会导致边界附近不够光滑（见 文献\cite{Lin2024PolynomialCC} 图 4）。
若输入二维形状仅由直边构成，则可由初始线性笼紧密逼近；这时，多项式 Cauchy 坐标可通过将直边变形为 \Bezier 曲线来获得更平滑结果，其依据是以下定理：
\begin{theorem}[文献\cite{Bell2015cauchytransform}的定理 3.3 ]\label{pro:cau-holomorphic}
设 \(\partial \Omega\) 光滑，且 \(u\) 是定义在 \(\Omega\) 上的全纯函数，并可延拓为 \(\Omega\cup \partial \Omega\) 上的连续函数。若 \(u\) 的边界值在 \(\partial\Omega\) 的某条开弧 \(\Gamma\) 上光滑，则 \(u\) 的各阶导数都可由 \(\Omega\) 连续延拓到 \(\Omega \cup \Gamma\)。
\end{theorem}

然而，对于曲边形状若仍使用线性输入笼，边界非全纯问题依旧存在（图~\ref{fig:cau-cmp-smoothness}）。
幸运的是，定理~\ref{pro:cau-holomorphic} 也表明：若变形前后某条边都光滑，则 $u$ 在该边界处可保证全纯。
这也推动我们将 Cauchy 坐标推广到高阶输入笼，使其能用光滑边更紧密包围输入形状。
假设 \(\partial \Omega\) 为分段光滑边界，可由定理~\ref{pro:cau-holomorphic} 直接推出：
 \begin{lemma}\label{pro:cau-lemma_holo}
 全纯函数 \(u\) 在 \(\partial \Omega\) 上全纯，当且仅当 \(f\) 将边界的光滑段映射到光滑段，且 \(f\) 在非光滑顶点处为相似变换。
 \end{lemma}

\paragraph{推导}
现在将 Cauchy 坐标扩展到由多项式曲线组成的高阶初始笼 $\partial \Omega$。
输入边 $\me_j$ 表示为次数 $m_j \geq 1$ 的 \Bezier 曲线：$z_{\me_j}(t) = \sum^{m_j}_{k=0} B_k^{m_j}(t) p^\text{old}_{\me_j,k} , \forall t \in [0, 1]$，其中 $p^\text{old}_{\me_j,k} \in \mathbb{C}$ 为控制点。
不失一般性，仍假设 $\me_j$ 的变形曲线为 $f^\text{new}(z_{\me_j}(t)) = \sum_{k=0}^{n_j} B^{n_j}_k(t) p^\text{new}_{\me_j,k}$。
则有
\begin{equation}\label{equ:cau-discrete-curved}
\begin{split}
    u(z) &=\frac{1}{2\pi i} \sum_{j=1}^n \int_{\me_j}\frac{f^\text{new}(\zeta)}{\zeta-z}d\zeta \\
    &= \frac{1}{2\pi i} \sum_{j=1}^n \int^1_{0} \frac{\sum_{k=0}^{n_j} p^\text{new}_{\me_j,k} B^{n_j}_k(t) z'_{\me_j}(t)}{z_{\me_j}(t) - z} dt \\
     &= \frac{1}{2\pi i} \sum_{j=1}^n \sum_{k=0}^{n_j} \left(\int^1_{0} \frac{ B^{n_j}_k(t) z'_{\me_j}(t)}{z_{\me_j}(t) - z} dt \right) p^\text{new}_{\me_j,k}, \ \ z \in \Omega.
\end{split}
\end{equation}
将式~\eqref{equ:cau-discrete-curved} 括号内积分记为 $\alpha_{\me_j,k}(z)$。
与控制点 $p^\text{new}_{\me_j,k}$ 相关的所有积分共同组成其 Cauchy 坐标。

\begin{equation}\label{equ:cau-derivative-curved}
     \frac{d^l \alpha_{\me_j,k}(z)}{ dz^l } = l! \int^1_{0} \frac{ B^{n_j}_k(t) z'_{\me_j}(t)}{(z_{\me_j}(t) - z)^{l+1}} dt.
\end{equation}
该式使坐标导数可直接计算。

\begin{figure}[t]
  \centering
  \begin{overpic}[width=0.99\linewidth]{chap4_smooth_poly}
    {
 \put(16,-2){\small \textbf{(a1)}}
 \put(40,-2){\small \textbf{(a2)}}
 \put(67,-2){\small \textbf{(b1)}}
 \put(89,-2){\small \textbf{(b2)}}
    }
  \end{overpic}
  \vspace{3mm}
  \caption{
  多项式 Cauchy 坐标将线性笼（a1，边上非端点控制点以实点表示）映射到三次笼时，边界附近会出现不光滑（a2）。
我们的方法使用的三次输入笼可紧密包围形状（b1），并得到光滑形变（b2）。
  }
  \label{fig:cau-cmp-smoothness}
\end{figure}

\subsection{解析解}\label{sec:cau-closed-form}
式~\eqref{equ:cau-discrete-curved}（坐标）和式~\eqref{equ:cau-derivative-curved}（坐标导数）中的被积函数均为有理多项式。
下面给出该类积分的解析解计算方法，这是本章方法的核心贡献。

\paragraph{表示方式}
将分母记为 $h(t), t\in[0,1]$，其次数为 $l$。
分子写为次数 $r$ 的多项式：$g(t) = \sum^r_{k=0} a_k t^k$，其中 $a_k$ 为系数。
由于 $\int^1_{0} g(t) / h(t) dt = \sum^r_{k=0} a_k \int^1_{0} t^k / h(t) dt$，只需计算 $\int^1_{0} t^k / h(t) dt, \forall k \geq 0$。

\paragraph{定理}
\begin{theorem}\label{pro:cau-general}
多项式 $h(t)$ 有 $m\leq l$ 个互异根 $\{\zeta_j\}^m_{j=1}$（$\zeta_j$ 为复数）。
则
$$\int_0^1 \frac{t^k}{ h(t)}dt=\sum_{j=1}^{m}\res(S_k(\zeta),\zeta_j), \zeta \in \mathbb{C}, $$
其中 $S_k(\zeta)=s_k(\zeta)/h(\zeta)$，且
$$s_k(\zeta)=\zeta^k(\log(1-\frac{1}{\zeta})+\sum_{j=1}^k\frac{1}{j \zeta^j}).$$
\end{theorem}

\begin{figure}[!t]
  \centering
  \begin{overpic}[width=0.5\linewidth]{chap4_cauchy_complex}
    {
 \put(63,19){\small $\zeta_i$}
 \put(52,70){\small $\zeta_j$}
 \put(65,54){\small $l_1$}
 \put(50,40){\small $l_2$}
 \put(27,36){\small $0$}
 \put(85,36){\small $1$}
 \put(85,58){\small $\gamma_2$}
 \put(27,58){\small $\gamma_1$}
 \put(84,80){\small $\Gamma$}
    }
  \end{overpic}
  \vspace{0mm}
  \caption{复平面上的积分路径。}
  \label{fig:cau-inte_path}
\end{figure}

\begin{proof}
下面考虑函数 \(S_k(\zeta)\) 沿积分路径（见图~\ref{fig:cau-inte_path}）的积分。
该路径由 \(\gamma_1,l_1,\gamma_2,l_2,\Gamma\) 组成，其中 \(\gamma_1,\gamma_2\) 分别是以 0 与 1 为圆心、半径为 \(\epsilon\) 的小圆，\(l_1,l_2\) 分别是线段 \([0,1]\) 的上、下侧，\(\Gamma\) 是半径为 \(R\) 的大圆（本证明令 \(R \to \infty\)）。

\(S_k(\zeta)\) 在该积分路径上全纯。
事实上，\(\log\left(1 - \frac{1}{\zeta}\right) = \log(1-\zeta) - \log(-\zeta)\)。
当 \(\zeta\) 沿任意绕过线段 \([0,1]\) 的简单闭曲线转一圈时，$\log\left(1 - \frac{1}{\zeta}\right)$ 仍保持单值全纯分支。
%
对 $S_k$ 在路径包围区域应用\emph{留数定理}，有
    \begin{equation}\label{equ:cau-residue_theorem}
        \begin{aligned}
 \int_{\gamma_1+l_1+\gamma_2+l_2+\Gamma} S_k(\zeta) \, d\zeta
            = 2\pi i \sum_{j=1}^{m} \res(S_k(\zeta), \zeta_j).
        \end{aligned}
    \end{equation}
当 $\zeta\in l_1$ 时，$\log(1-\zeta)-\log(-\zeta)=\log(1-t)+2 \pi i-\log(-t)$；
当 $\zeta\in l_2$ 时，$\log(1-\zeta)-\log(-\zeta)=\log(1-t)-\log(-t)$，其中 $t$ 为 $\zeta$ 的实部。
于是
\begin{equation}\label{equ:cau-l1l2}
\begin{aligned}
    &\int\limits_{l_1+l_2} S_k(\zeta) \, d\zeta=\int_{0}^1 \left(S_k(t) +\frac{2\pi i t^k}{h(t)}\right)\, dt+\int_1^0 S_k(t) \, dt\\
    &=2\pi i \int_{0}^{1} \frac{ t^k}{h(t)} \, dt.
    \end{aligned}
\end{equation}
 当 \( \zeta \in \Gamma \) 时，根据泰勒展开，\( s_k \) 可写为
    \begin{equation}
        s_k(\zeta)=\zeta^k(\sum_{j=1}^{\infty}\frac{-1}{j \zeta^j}+\sum_{j=1}^k\frac{1}{j \zeta^j})=\sum_{j=1}^\infty\frac{-1}{(k+j) \zeta^{j}}.
    \end{equation}
    由于 $h(\zeta)$ 次数至少为 1，有 $\lim_{R\to \infty}\zeta S_k(\zeta)=0$，从而
    \begin{equation}\label{equ:cau-Gamma}
        \lim\limits_{R\to \infty}\int_{\Gamma} S_k(\zeta) \, d\zeta=0.
    \end{equation}
    当 $\zeta \in \gamma_1$ 时，有
    \begin{equation}
    \begin{aligned}
       \lim\limits_{\epsilon\to 0} &\int_{\gamma_1} S_k(\zeta) \, d\zeta=\lim\limits_{\epsilon\to 0}\int_{\gamma_1} \frac{\zeta^k \log(-\zeta)}{h(\zeta)}\, d\zeta.\\
       &\leq \frac{\epsilon^k (-\log(\epsilon)+2 \pi)}{h(0)}2 \pi \epsilon \to 0.
       \end{aligned}
    \end{equation}
    对 $\zeta \in \gamma_2$ 同理，因此
    \begin{equation}\label{equ:cau-gamma12}
       \lim\limits_{\epsilon\to 0} \int_{\gamma_1+\gamma_2} S_k(\zeta) \, d\zeta=0.
    \end{equation}

    将~\eqref{equ:cau-l1l2}、\eqref{equ:cau-Gamma}、\eqref{equ:cau-gamma12} 代入~\eqref{equ:cau-residue_theorem}，即可得到
    \begin{equation}
        \int_0^1 \frac{t^k}{ h(t)}dx=\sum_{j=1}^{m}\res(S_k(\zeta),\zeta_j).
    \end{equation}
\end{proof}

\paragraph{说明}
在上述证明中，$\log(1-\zeta) - \log(0-\zeta)$ 将积分区间限定为 \([0,1]\)，额外项 $\sum \frac{1}{j \zeta^j}$ 用于保证 \(S_k\) 在无穷远点的留数为 0。
该定理也可扩展到区间 \([a,b]\) 的积分，只需把 $s_k(\zeta)$ 中对数项改写为 \(\log(b-\zeta)-\log(a-\zeta)\)，并保持附加项为该对数项泰勒展开前 $k$ 项的相反数。

\paragraph{一般根}
在最一般情形下，\( \zeta_j \) 是 \( S_k(\zeta) \) 的单极点。记 $h^{j,1}(\zeta)=h(\zeta)/(\zeta-\zeta_j)$，则
\begin{equation}\label{equ:cau-gene_expression}
\begin{aligned}
    &\res(S_k(\zeta),\zeta_j)=\lim_{\zeta\to \zeta_j}(\zeta-\zeta_j)S_k(\zeta)=\frac{s_k(\zeta_j)}{h^{j,1}(\zeta_j)}.
    \end{aligned}
\end{equation}

\paragraph{重根}
在计算 \(\alpha_{\me_j,k}(z)\) 时，少量点会出现多重极点。
例如，当输入边 $\me_j$ 为二次曲线时，多重极点意味着 \(z\) 与 $\me_j$ 的焦点重合。
此时记 \( \zeta_j \) 为 $l_j$ 阶极点，且 $h^{j,l_j}(\zeta)=h(\zeta)/(\zeta-\zeta_j)^{l_j}$，有
\begin{equation}\label{equ:cau-repeated_expression}
\begin{aligned}
&\res(S_k(\zeta),\zeta_j)
    =\lim_{\zeta\to \zeta_j}\frac{1}{(l_j-1)!}\frac{d^{l_j-1}}{d\zeta^{l_j-1}}\{\frac{s_k(\zeta)}{h^{j,l_j}(\zeta)}\}.
    \end{aligned}
\end{equation}
对给定 $l_j$，按标准求导规则即可计算。
在实际计算中，尚未遇到计算 \(\alpha_{\me_j,k}(z)\) 时 $l_j>1$ 的情形。

\subsection{扩展}\label{sec:cau-connection}
\paragraph{多项式 Green 坐标}
文献\cite{MichelThiery2023} 提出了闭式多项式 Green 坐标，可将直边输入笼变形为曲边输出笼；但其不能处理高阶输入笼。
在此基础上，我们利用上述定理将 Green 坐标扩展到高阶笼。
根据文献~\cite{MichelThiery2023} 的式 (12)，高阶笼 Green 坐标计算可归结为如下积分：
定理~\ref{pro:cau-general} 可直接用于该积分。

\paragraph{有理多项式笼}
我们还可将坐标扩展到由有理多项式曲线构成的初始高阶笼 $\partial \Omega$（图~\ref{fig:cau-rational}）。
令 $z_{\me_j}(t) = \frac{z_{w,\me_j}(t)}{w(t)}, \forall t \in [0, 1]$，其中 \(z_{w,\me_j}(t)=\sum_{k=0}^{m_j}w_{\me_j,k}p_{\me_j,k}^{\text{old}}B_k^{n_j}(t)\)，\(w(t)=\sum_{k=0}^{m_j}w_{\me_j,k}B_k^{n_j}(t)\)。
随后将 $f^\text{new}(z_{\me_j}(t))$ 写为与其共享权重、控制点为新控制点 \(p^\text{new}_{\me_j,k}\) 的有理曲线，可得
\begin{equation}
    \alpha_{\me_j,k}(z)=\int_0^1 \frac{z_{w,\me_j}'(t)w(t)-w'(t)z_{w,\me_j}(t)}{(z_{w,\me_j}(t)-zw(t))w(t)^2}dt.
\end{equation}
定理~\ref{pro:cau-general} 同样可用于该积分。
当 \(w(t)\) 与 \(z_{\me_j}^{w}(t)-zw(t)\) 的根重合时，需要额外处理：实践中以误差阈值 \(\epsilon=10^{-10}\) 合并根后重新评估其重数。

\subsection{实现细节}\label{sec:cau-details}
算法~\ref{alg:cau-encoding} 给出了本章 Cauchy 坐标的计算伪代码。

\paragraph{计算 \(\{\zeta_j\}\)}
为应用定理~\ref{pro:cau-general} 计算高阶笼 Cauchy 坐标，需要先求解 $h(t)=0$ 的根 $\{\zeta_j\}$。
对最常用的二次和三次曲线，可使用闭式求根公式。
对四次以上多项式，则需数值求根。

\paragraph{递推计算}\label{sec:cau-encoding}
以下以迭代方式计算 \(\res(S_{k}(\zeta),\zeta_j)\)。为简洁起见，仅给出 \(\zeta_j\) 为单极点的情形（多极点类似）。
由
    $s_{k+1}(\zeta)=\zeta s_k(\zeta)+\frac{1}{k+1},\ k\geq 0,$
在实践中使用如下递推：
\begin{equation}\label{equ:cau-recur}
\begin{aligned}
&\res(S_{0}(\zeta),\zeta_j)=\frac{s_0(\zeta_j)}{h^{j,1}(\zeta_j)},\\
    &\res(S_{k+1}(\zeta),\zeta_j)=\zeta_j \res(S_k(\zeta),\zeta_j)+\frac{1}{k+1}\frac{1}{h^{j,1}(\zeta_j)}.
    \end{aligned}
\end{equation}

\begin{algorithm}[t]
  \caption{计算初始位置点 $z$ 的 Cauchy 坐标}\label{alg:cau-encoding}
  \SetKwInput{KwInput}{输入}
  \SetKwInput{KwOutput}{输出}
  \SetKwFunction{FMain}{Main}
  \SetKwFunction{FD}{D}

  \KwInput{初始笼内部一点 $z$，以及由次数 $m_j$ 多项式 $z_{\me_j}(t)$ 表示的初始边 $\me_j$。}
  \KwOutput{点 $z$ 关于变形边控制点 \(p_r^{\text{new}}\) 的坐标 $\alpha_{\me_j,r}(z)$。}

  计算 $z_{\me_j}(t)-z=0$ 的根 $\{\zeta_j\}_{j=1}^l$\;
  将 \(B_k^{n_j}(t)z_{\me_j}'(t)\) 转换为 \(\sum_{k=0}^{m_j+n_j-1} c_k t^k\)\;
  $k \gets 0; \alpha_{\me_j,r}(z) \gets 0$\;
  \For{$k \leq m_j+n_j-1$}{
  $temp \gets 0$\;
      \For{每个 \(\zeta_j\) 属于 \(\{\zeta_j\}\)}{
      \If{\(\zeta_j\) 为单极点}{用 \eqref{equ:cau-recur} 计算 \(\res(S_{k}(\zeta),\zeta_j)\)。}
      \Else{用 \eqref{equ:cau-repeated_expression} 计算 \(\res(S_{k}(\zeta),\zeta_j)\)。}
      $temp \gets temp+\res(S_k(\zeta),\zeta_j))$;
      }
      \(\alpha_{\me_j,r}(z) \gets \alpha_{\me_j,r}(z) + c_k temp; \, k \gets k+1 \)\;
  }
\end{algorithm}

\paragraph{数值稳定性}
在使用~\eqref{equ:cau-recur} 计算时，只需显式计算 $1/h^{j,1}(\zeta_j)$ 与 $s_0(\zeta_j)$。
当 \(\zeta_j\) 是单根时，$1/h^{j,1}(\zeta_j)$ 定义良好。
函数 \( s_0(\zeta) \) 可写为：
{\small
\[
\log(1-\frac{1}{\zeta}) = \frac{1}{2} \log\left(1+\frac{ 1 - 2 \operatorname{Re}(\zeta)}{\|\zeta\|^2}\right) + i \text{atan2}\left(\operatorname{Im}(\zeta), \|\zeta\|^2 - \operatorname{Re}(\zeta)\right).
\]
}
当 \( z \notin \partial \Omega \) 时该表达是良定义的：
(1) 对数函数 \( \log(\cdot) \) 仅在 \( \|\zeta\|^2 = 0 \) 或 \( 1 - 2 \text{Re}(\zeta) \leq -\|\zeta\|^2 \) 时失效，这等价于 \( \zeta = 0 \) 或 \( \zeta = 1 \)（即 \( z \) 在 $\me_j$ 两端点上）；
(2) 反正切函数在 \( \zeta \neq 0 \) 或 \( \zeta \neq 1 \) 时同样定义良好。
我们对坐标进行数值误差检验，在所有样例中最大误差未超过 $10^{-11}$。

\section{实验与评估}\label{sec:cau-results}
本章方法采用 C++ 实现。
所有实验均在一台台式机上完成，配置为 4.00 GHz Intel Core i7-4790K 处理器和 16 GB 内存。
%
除图~\ref{fig:cau-rational} 外，其他图中的高阶输入笼边均为多项式曲线。

\paragraph{基于笼的变形}
本章在图~\ref{fig:cau-more-example} 中展示了多组图像变形示例。
本章坐标能够对有机形状和机械形状提供光滑且保角的变形结果。
本章方法也支持对曲线段多项式次数进行升阶和降阶（图~\ref{fig:cau-diff-order}）。
%\paragraph{Different orders}
%We allow \( n_j < m_j \); 
当输出次数低于输入次数时，先对曲线降阶会改变笼的初始形状，但不会影响艺术家后续的编辑过程。
此外，引理~\ref{pro:cau-lemma_holo} 保证了所得变形在边界附近仍保持光滑。
\begin{figure}[t]
	\centering
	\begin{overpic}[width=0.99\linewidth]{figures/chap5_cau_diff_order_new.pdf}
		{
		}
	\end{overpic}
	\vspace{-2mm}
	\caption{
		高阶输入笼可被变形为不同次数的高阶输出笼。上排输入笼为二次曲线，下排输入笼为三次曲线。
	}
	\label{fig:cau-diff-order}
\end{figure}

\paragraph{点到点（P2P）变形}
%To achieve better control over the cage's deformation, 
为支持由把手驱动的变形，我们利用坐标导数构建目标函数并进行优化求解，这与文献 \cite{Weber2009} 和 文献\cite{Lin2024PolynomialCC} 的做法类似。
目标函数包含把手点之间距离平方和，以及一个由边界上坐标二阶导数的 $\ell_2$ 范数组成的光滑项。
%
该优化问题是最小二乘问题。
通过线性求解得到最优解。
此外，在线性系统中，拖动把手时系数矩阵保持不变；因此我们在预处理阶段仅做一次预分解，从而支持交互式把手编辑（图~\ref{fig:cau-teaser}、\ref{fig:cau-more-example} 和 \ref{fig:cau-diff-order}）。
%Minimizing this function is a least-squares problem, allowing users to edit handles interactively (Fig. \ref{fig:cau-teaser}, \ref{fig:cau-diff-order}).
%for real-time handle editing by the user (Fig. \ref{fig:cau-teaser}, \ref{fig:cau-diff-order}).}
%point-to-point term 


\begin{figure}[t]
	\centering
	\begin{overpic}[width=0.99\linewidth]{figures/chap4_comparep2p.pdf}
		{
			\put(6,0){\small 线性输入笼}
			\put(56,0){\small 高阶输入笼}
			\put(29,0){\small Cauchy 坐标}
			\put(86,0){\small 本章方法}
		}
	\end{overpic}
	\vspace{2mm}
	\caption{
		在相同控制点数量下，与经典 Cauchy 坐标的点到点变形对比。红框突出显示了边界附近的光滑性差异。
	}
	\label{fig:cau-comparep2p}
\end{figure}

\begin{figure}[t]
	\centering
	\begin{overpic}[width=0.99\linewidth]{figures/chap4_comparep2p_poly.pdf}
		{
			\put(3,-2){\small 多项式 Cauchy 坐标}
			\put(30,-2){\small 本章方法}
			\put(45,-2){\small 多项式 Cauchy 坐标}
			\put(79,-2){\small 本章方法}
		}
	\end{overpic}
	\vspace{2mm}
	\caption{
		与多项式 Cauchy 坐标的点到点变形对比。两种方法的输入笼具有相同数量的控制点。
	}
	\label{fig:cau-comparep2p_poly}
\end{figure}

\paragraph{时间与内存}
%All the cage-based deformation algorithms run in real time for all the examples shown in this paper. 
%All cage-based and P2P deformations achieve real-time performance across all demonstrated examples. 
在文中展示的所有案例中，基于笼与 P2P 两类变形均可实时运行。
%We focus solely on analysing the pre-computation time. 
本章只分析预计算时间。
如第~\ref{sec:cau-details} 节所述，一旦根 \(\{\zeta_j\}\) 已确定，积分解析式的求值只涉及少量复数乘法和对数运算。
在实践中，使用闭式表达计算 \(\{\zeta_j\}\) 会占据较多预计算时间（对二次曲线约占 24\%，对三次曲线约占 33\%）。
%
%During precomputation, our coordinates are based on the integral convolutions of monomials, allowing artists to precompute the coordinates up to a certain degree \( n_j \). This flexibility lets them adjust the degree of the editing curves \(  \leq n_j \) at runtime without the need for coordinate recomputation. 
因此，在 \(\{\zeta_j\}\) 计算完成后，先把有理多项式积分预计算到一个较高指定次数会更方便。
%This approach gives artists the runtime flexibility to modify the degree of editing curves to any value less than or equal to the specified degree, without necessitating coordinate recomputation.
%
这种做法使艺术家能够在不重算坐标的前提下，自由调整曲线编辑次数（不超过预设上限）。
% In practice, once the degree of the deformed curves is determined, we express the control points of the polynomial basis as a linear combination of the control points of the Bernstein basis. When subjected to this linear transformation, our coordinates yield the coordinates of the control points of the \Bezier curve.
%the total memory during preprocessing does not exceed 90 MB
在测试中，预处理内存占用不超过 90 MB。
该上限由图~\ref{fig:cau-cmp-cages} 中的 Crab 模型达到：其笼由 24 条三次曲线组成，内部包围约 27000 个点。

\paragraph{与柯西坐标~\cite{Weber2009} 的比较}
Cauchy 坐标~\cite{Weber2009} 将线性笼映射为线性笼。
%\todo{closeness}Cauchy coordinates
线性笼通常需要过高自由度才能表示复杂曲边形状，这会显著削弱交互操控性（图~\ref{fig:cau-cmp-cages} 与 \ref{fig:cau-cmp-linearcages}）。
%In contrast, the curved cage achieves more natural deformations with fewer control points, enabling intuitive shape editing through fewer adjustments.
相比之下，曲边笼用更少控制量即可产生更自然的变形，且只需更少调整即可完成编辑。
%
在图~\ref{fig:cau-cmp-cages} 的示例中，线性笼方法需要更复杂操作和更长时间，才能接近高阶笼结果。
%achieve deformation comparable to high-order-cage results. % and~\ref{fig:cau-cmp-linearcages}
%
此外，线性笼与复杂形状之间常存在空隙，导致笼附近变形不光滑（见图~\ref{fig:cau-rational} 的基于笼变形和图~\ref{fig:cau-comparep2p} 的 P2P 变形）。
而高阶笼具有更高几何贴合度，因此能得到更光滑的变形。
%\todo{smoothness}


\begin{figure}[t]
  \centering
  \begin{overpic}[width=0.95\linewidth]{chap4_zoomin_mvc} %
    {
    \put(7,-0.5){\small \textbf{(a)}}
 \put(23,-0.5){\small \textbf{(b)}}
 \put(43,-0.5){\small \textbf{(c1)}}
 \put(63,-0.5){\small \textbf{(c2)}}
 \put(85,-0.5){\small \textbf{(c3)}}
    }
  \end{overpic}
  \vspace{3mm}
  \caption{
  (a) 输入线性笼。
  (b) 三次 MVC。
(c1, c2, c3) 对输入边 (a) 采用三种不同初始参数化时，采用我们的方法得到的变形结果。
  %Starting from a linear cage, we apply different parameterizations to one of its edges to obtain various deformation results (c1, c2, c3). Compared to MVC (b), our coordinates achieve conformal deformation, preserving local details of the image (a).
  }
  \label{fig:cau-cmp-mvc}
\end{figure}

\begin{figure}[t]
  \centering
  \begin{overpic}[width=0.95\linewidth]{chap4_compare2polyGC} %
    {
    \put(6.5,6.5){\small 初始姿态}
 \put(53,6.5){\small 初始姿态}
    \put(18,-2){\small 多项式 Cauchy 坐标}
\put(78,-2){\small 我们的方法}
    }
  \end{overpic}
  \vspace{2mm}
  \caption{
  与多项式 Cauchy 坐标比较：输入线性笼（左）和高阶笼（右）控制点数量相同，且均变形为三次笼。
  %, our cage remains closer to the image both initially and after deformation, resulting in more intuitive deformations.
  }
  \label{fig:cau-cmp-pgc}
\end{figure}




\begin{figure}[t]
	\centering
	\begin{overpic}[width=0.92\linewidth]{figures/chap4_rational.pdf}
		{
			\put(10,-1){\small \textbf{(a1)}}
			\put(35,-1){\small \textbf{(a2)}}
			\put(63,-1){\small \textbf{(b1)}}
			\put(90,-1){\small \textbf{(b2)}}
		}
	\end{overpic}
	\vspace{3mm}
	\caption{
		在相同控制点数量下，线性笼 (a1) 与二次有理笼 (b1) 的对比。采用高阶输入笼得到的变形结果 (b2) 比线性笼上的 Cauchy 变形 (a2) 更光滑。
	}
	\label{fig:cau-rational}
\end{figure}

\begin{figure}[t]
	\centering
	\begin{overpic}[width=0.99\linewidth]{figures/chap5_cau_gallery.pdf}
		{
		}
	\end{overpic}
	\vspace{0mm}
	\caption{
		利用本章坐标及其导数在高阶输入笼上实现的多组基于笼变形与点到点变形结果。
	}
	\label{fig:cau-more-example}
\end{figure}



\begin{figure}[t]
	\centering
	\begin{overpic}[width=0.99\linewidth]{figures/chap4_cage_complex.pdf}
		{
			\put(1,-1.5){\small 线性输入笼（151 点）}
			\put(43,-1.5){\small 三次输入笼（41 点）}
			\put(22,-1.5){\small Cauchy 坐标的 P2P}
			\put(64,-1.5){\small 本章坐标的 P2P}
			\put(82,-1.5){\small 本章笼变形}
		}
	\end{overpic}
    \vspace{2mm}
	\caption{
		线性输入笼与高阶输入笼的对比。线性笼需要大量边段才能包围给定形状，而三次高阶笼只需更少参数即可完成包围与编辑。
	}
	\label{fig:cau-cmp-linearcages}
\end{figure}


\paragraph{与三次均值坐标~\cite{Li2013} 的比较}
三次均值坐标（Cubic MVCs）~\cite{Li2013} 支持将直边笼变形为三次曲线。
%
由于线性笼是高阶笼的特例，本章坐标也可用于直边笼并实现保角变形。
%
我们通过为线性输入笼的每条边指定不同的非线性参数化，获得不同的保角变形结果（图~\ref{fig:cau-cmp-mvc}）。
但 Cubic MVC 的结果存在较高畸变。

\paragraph{与多项式柯西坐标~\cite{Lin2024PolynomialCC} 的比较}
%polynomial Cauchy coordinates, polynomial Green coordinate
在给定相同线性输入笼与高阶变形笼时，多项式 Cauchy 坐标~\cite{Lin2024PolynomialCC} 与多项式 Green 坐标~\cite{MichelThiery2023} 产生相同变形。
因此，这里只比较使用直边输入笼时的多项式 Cauchy 坐标~\cite{Lin2024PolynomialCC}。
%
为公平比较，多项式 Cauchy 坐标与本章方法使用相同数量的控制点。
在其变形中，变形曲线控制点在初始姿态共线，因此线性输入笼可能缺乏足够自由度来紧贴输入形状。
%
%For initially straight shapes, we compare our method with the polynomial Green coordinate (PolyGC)~\cite{MichelThiery2023}. 
%Due to its conformal deformation results, PolyGC is often preferred by artists over Cubic MVC~\cite{Li2013} in shape deformation applications. 
当线性输入笼无法贴合复杂形状边界时，变形后边界附近会出现不光滑（图~\ref{fig:cau-cmp-smoothness}）；并且相较其方法，采用我们的方法得到的变形形状更贴近变形笼（图~\ref{fig:cau-cmp-pgc}）。
%closeness
%Moreover, %given the same number of control points, higher-order cages approximate the shape's contour more closely than linear cages. %This means that artists can achieve editing results that better align with their intentions when manipulating the cage.
%as shown in Fig.~\ref{fig:cau-cmp-pgc}, the deformed shape is closer to the deformed cage using the higher-order cage via our coordinates than theirs.
%This enables artists to obtain deformations that match their creative intent more precisely through high-order cage-based manipulation.
因此，艺术家通过高阶笼操控可获得更光滑、也更符合意图的变形。
%
最后，在图~\ref{fig:cau-comparep2p_poly} 中，我们比较了在输入笼自由度相同条件下的 P2P 变形。
%
在这些案例里，其线性笼无法紧密贴合输入形状（例如图~\ref{fig:cau-comparep2p_poly} 中大象后腿之间没有分离），从而产生不自然的 P2P 变形。
高阶笼能够紧密包围形状，可避免该问题。

\begin{figure}[t]
  \centering
  \begin{overpic}[width=0.965\linewidth]{chap4_compare_cauchy} %
    {
    % \put(17,25){\small (a1)}
    % \put(44,25){\small (a2)}
    % \put(30,-2){\small (c)}
    % \put(85,25){\small (b)}
    % \put(85,-2){\small (d)}
     \put(2,24){\small 高阶输入笼}
    \put(36,24){\small 中间笼}
    \put(18,-2){\small 高阶输入笼}
    \put(62,25.5){\small 多项式 Cauchy 坐标}
\put(84,-2){\small 我们的方法}
    }
  \end{overpic}
  \vspace{3mm}
  \caption{
在高阶输入笼上的比较。其方法需要两个输入笼、需要后验调整且不具备闭式表达（上）；我们的方法只需一个输入笼，可直接变形且具有闭式表达（下）。 %does not require post-adjustments
  }
  \vspace{-2mm}
  \label{fig:cau-cmp-poly-Cauchy}
\end{figure}

%\todo{high-order cages Fig.~\ref{fig:cau-cmp-poly-Cauchy}}
%PolyGC can not be used for the high-order input cages.
%For initially highly curved shapes, PolyGC may fail to achieve the desired edits when restricted to the same limited number of control points as our method. Therefore, 
%\tr{For such cases, we compare our approach against GC \cite{Lipman2008} and polynomial Cauchy coordinates \cite{Lin2024PolynomialCC}. Compared to GC in linear cages, the superior fitting capability of high-order cages enables artists to achieve more intuitive and flexible deformations (Fig. \ref{fig:cau-cmp-cages}). 
%\todo{double check!!!}
多项式 Cauchy 坐标可用于高阶输入笼。
具体而言，其做法是在中间笼上应用多项式 Cauchy 坐标，以生成用于包围输入形状的高阶输入笼（图~\ref{fig:cau-cmp-poly-Cauchy}）。
随后再通过数值方法计算从中间笼到输入笼映射的逆映射，进而得到形状变形。
%
本章坐标具有以下优势。
首先，由于多项式 Cauchy 坐标不具插值性，其高阶输入笼生成方法无法产生任意有效的高阶输入笼。
%their method cannot compute valid coordinates on arbitrarily shaped curved input cages. % due to their non-intuitive geometric constraints on the high-order input cages. %, significantly limiting its applicability to complex geometries.
%the cage enclosing the shape can be directly manipulated via our coordinates  without
%Second, they require subsequent post-adjustments to the high-order input cages according to the intermediate cages, whereas we do not need it%our coordinates enable direct manipulation without adjustments (Fig.~\ref{fig:cau-cmp-poly-Cauchy}).
%Second, %they need post-adjustments to high-order cages based on intermediate cages, while we do not .
对于含狭窄区域的输入形状（如图~\ref{fig:cau-comparep2p_poly} 中章鱼触爪之间或大象后腿之间区域），其高阶输入笼很难调整到既避免自交又保持双射，从而导致算法失败。
此外，要让其高阶输入笼精确贴合形状几乎不可能。
%it is difficult to make their high-order input cages avoid self-intersections, losing bijectivity to cause their algorithm failure.
%due to the non-interpolating coordinates, the deformation from the intermediate cage to the initial cage can result in self-intersections. Such self-intersections signify a loss of invertibility, ultimately causing their algorithm to fail.}
%
%Second, since they use numerical methods for the inverse of the mapping from the intermediate to the input cages, their coordinates have larger errors and long computation time (see Table 1 in~\cite{Lin2024PolynomialCC}) and cannot be used for P2P deformation due to the lack of closed forms. 
第二，其坐标不具闭式表达，因此计算误差更大且耗时更长（见~\cite{Lin2024PolynomialCC} 的表 1），并且无法用于 P2P 变形。

%\tr{Compared with polynomial Cauchy coordinates, our method offers the advantage that  Furthermore, because they use numerical methods for inverse mapping computation, their resulting coordinates have larger errors and longer computation times (see their Table 1). Finally, due to the lack of an explicit expression for the coordinates, they are unable to implement point-to-point (P2P) deformation within high-order cages.}



% We select the Cubic MVC~\cite{Li2013} and the  Polynomial Green Coordinate (PolyGC)~\cite{MichelThiery2023} as the competitors.
% %
% Since Cubic MVC makes the output cages three-order, we set the order of the output polynomial curves as three for PolyGC and our coordinate in comparison.
% %
% Since our coordinate is defined on a high-order cage, 
% %Different from PolyGC, 
% we elevate the orders of the input cage's edges to three if the input order is less than three.
% Specifically, we apply the non-uniform parameterization %\tr{express the straight edge as a cubic polynomial curve} 
% during this elevation, resulting in a non-uniform distribution of control points (e.g., Figure~\ref{fig:cau-cmp} (c1) and (c2)).


% We compare with Cubic MVC in Figures~\ref{fig:cau-cmp-mvc} and~\ref{fig:cau-cmp}.
% Cubic MVC is built upon the Green identity on the unit disk, which helps diffuse both Dirichlet and Neumann conditions on the unit projection sphere. Cubic MVC extends the capability to interpolate boundary conditions and define the deformation function across the entire space, not just within the cage. Cubic MVC retains the interpolation property, whereas we trade interpolation for conformality. 
% The results corresponding to our various parameterizations consistently ensure that the local aspects of details are preserved precisely, as shown in Figure~\ref{fig:cau-cmp-mvc}.


% %\paragraph{Comparison with polynomial GC}
% We compare with PolyGC in Figures~\ref{fig:cau-cmp-pgc} and~\ref{fig:cau-cmp}. 
% %
% In Figure~\ref{fig:cau-cmp-pgc} and the bottom row of Figures~\ref{fig:cau-cmp}, we construct the rest polygonal cages and higher-order cages that are close to the shape. Both cages have the same number of edges and the same number of control points during deformation. The cage is then deformed to be high-order with the same order. % resulting in the final deformed shape. 
% Both our coordinate and PolyGC enable smooth editing while ensuring conformal deformation results. Compared to PolyGC, our coordinates produce more intuitive deformations, as our cage remains closer to the shape both in the initial and deformed states.
% % \tr{We note that when the rest shape is highly curved as in Figure~\ref{fig:cau-cmp-cages}, pGC cannot achieve the desired deformation with the same number of control points as ours.
% % }
% In the first and second rows of Figure~\ref{fig:cau-cmp-pgc}, we use the same polygonal cage but with different initial parameterizations on some edges (i.e., the positions of the control points are different). 
% In this case, our method can be seen as a complement to PolyGC. Applying different high-order parameterizations to straight edges can produce varying deformation results.
%  %In such cases, applying different higher-order parameterizations to straight edges can produce varying deformation results (see Figure~\ref{fig:cau-cmp-mvc} and the top and middle of Figure~\ref{fig:cau-cmp}). 
% When the parameterization is concentrated on one side of the segment, that side becomes more sensitive to the shape transformation of the cage boundary, making the deformation closer to the cage's shape (see the top of the pants and the tail of the fish in Figure~\ref{fig:cau-cmp}).



\section{本章小结}\label{sec:chap5-cauchy-conclusion}
本章提出了一种定义在高阶输入笼上的二维 Cauchy 坐标及其导数的闭式解析方法。为实现这一目标，我们从 Cauchy 变换出发，将由多项式曲线组成的高阶笼纳入统一的坐标框架，并将坐标计算转化为有理多项式积分的解析求值。基于对数函数与留数定理，我们给出了坐标及其导数的闭式表达，同时说明了同一积分框架向高阶 Green 坐标扩展的方式。实验结果表明，本章方法能够支持基于笼的交互式变形与点到点变形，在保持保角特性的同时获得更光滑、更贴合边界且更符合编辑意图的变形结果；与线性笼及已有相关方法相比，本章方法在高阶输入笼场景下具有更好的几何适应性与实用性。
